\documentclass[11pt,letterpaper]{article}
\usepackage[margin=1in]{geometry}
\usepackage{amsmath,amssymb}
\usepackage{graphicx}
\usepackage{booktabs}
\usepackage{xcolor}
\usepackage{float}
\usepackage{colortbl}
\usepackage{titlesec}
\usepackage{enumitem}
\usepackage{multirow}
\usepackage{hyperref}

\definecolor{phasebg}{HTML}{E3F2FD}
\definecolor{resultbg}{HTML}{E8F5E9}
\definecolor{newbg}{HTML}{FFF9C4}
\definecolor{seccolor}{HTML}{1565C0}

\titleformat{\section}{\large\bfseries\color{seccolor}}{\thesection}{0.6em}{}[\titlerule]
\titleformat{\subsection}{\normalsize\bfseries}{\thesubsection}{0.5em}{}

\title{\textbf{Project Progress Report}\\[0.3em]
  \large Classification and Measurement of One-Dimensional Hyperuniform Point Patterns}
\author{Michael Fang \quad --- \quad Advisor: Prof.\ Salvatore Torquato, Princeton University}
\date{March 27, 2026}

\begin{document}
\maketitle
\tableofcontents
\vspace{1em}
\noindent\rule{\linewidth}{0.4pt}

%==============================================================================
\section{Project Overview}
%==============================================================================

This project provides a systematic numerical characterization of hyperuniform point
patterns in one dimension. Two complementary metrics are computed for each pattern:
\begin{itemize}[noitemsep]
  \item \textbf{Hyperuniformity exponent} $\alpha$: characterizes how $S(k) \sim |k|^\alpha$
        as $k \to 0$.
  \item \textbf{Surface-area coefficient} $\bar\Lambda$: measures residual disorder within
        the strongest (Class~I) patterns; $\bar\Lambda = 1/6$ is the global minimum (integer lattice).
\end{itemize}

\noindent Three universality classes arise from the large-$R$ behavior of the number
variance $\sigma^2(R)$:

\vspace{0.4em}
\noindent\begin{tabular}{llll}
\toprule
\textbf{Class} & \textbf{Condition} & \textbf{Variance} & \textbf{$\bar\Lambda$} \\
\midrule
I   & $\alpha > 1$ & bounded                 & finite \\
II  & $\alpha = 1$ & $\sim \Lambda_{II} \ln R$     & diverges \\
III & $0 < \alpha < 1$ & $\sim \Lambda_{III} R^{1-\alpha}$ & diverges \\
\bottomrule
\end{tabular}
\vspace{0.5em}

\noindent The project is organized by ``weeks,'' each corresponding to a checkpoint meeting
with Prof.\ Torquato (not necessarily one calendar week). Three such meetings have
taken place to date, yielding 27 characterized patterns and a Junior Paper
(§1--§3 complete, 9 pages; results sections pending). A fourth meeting cycle
is now underway.

%==============================================================================
\section{Literature Review}
%==============================================================================

\noindent This section surveys the primary literature that motivates, informs, and
contextualizes the project. Twelve papers are reviewed, organized thematically.

%------------------------------------------------------------------------------
\subsection{Foundations of Hyperuniformity}
%------------------------------------------------------------------------------

\noindent\textbf{Torquato \& Stillinger (2003).}
The foundational paper defines \textbf{hyperuniformity} as the condition
$S(\mathbf{k}\to\mathbf{0}) = 0$, where $S(\mathbf{k})$ is the structure factor.
For a $d$-dimensional point pattern the \textbf{number variance}
$\sigma^2(R) \equiv \langle N(R)^2\rangle - \langle N(R)\rangle^2$
grows at most as the surface area $R^{d-1}$ for large $R$, in contrast to the
$R^d$ growth of a Poisson process.
Three universality classes are introduced: Class~I ($\alpha>1$, bounded $\bar\Lambda$),
Class~II ($\alpha=1$, $\sigma^2\sim \Lambda_{II}\ln R$), Class~III ($0<\alpha<1$,
$\sigma^2\sim \Lambda_{III}R^{1-\alpha}$).

The \textbf{surface-area coefficient}
\[
  \bar\Lambda = \lim_{R\to\infty} \frac{\sigma^2(R)}{2\rho R^{d-1}A_1(R)}
\]
(with $A_1(R)$ the window surface area) serves as an order metric for Class~I patterns.
For $d=1$ the integer lattice achieves
$\bar\Lambda_{\min} = 1/6$, derived from the exact formula
$\sigma^2(R) = \{R\}(1-2\{R\})$ where $\{R\}$ is the fractional part of $R$.
In higher dimensions, Bravais lattices minimize $\bar\Lambda$ among periodic patterns.
The paper establishes the conceptual framework adopted throughout this project.

\vspace{0.4em}
\noindent\textbf{Zachary \& Torquato (2009).}
This paper extends hyperuniformity to \textbf{two-phase heterogeneous media}.
In this setting the relevant quantity is the \textbf{spectral density}
$\tilde\chi_V(\mathbf{k})$ --- the Fourier transform of the two-point correlation
of the indicator function $I(\mathbf{x})$ for phase 2.
Hyperuniformity now requires $\tilde\chi_V(\mathbf{k}\to\mathbf{0}) = 0$, in direct
analogy with $S(\mathbf{k}\to\mathbf{0}) = 0$ for point patterns.

A central technical contribution is the \textbf{theta-function (Bragg-sum) method}
for computing $\bar\Lambda$ analytically from the Fourier-space structure of the
pattern: Bragg peaks at reciprocal-lattice vectors $\mathbf{k}_j$ contribute to
$\sigma^2(R)$ through overlap integrals of the window function.
This method is applied to tabulate $\bar\Lambda$ for $d=1,2,3$ Bravais lattices and
one-dimensional quasicrystals.
The Fibonacci chain result \textbf{$\bar\Lambda = 0.20110$} (Table~1 of that paper)
serves as the primary literature benchmark for this project's numerical validation.

%------------------------------------------------------------------------------
\subsection{Substitution Tilings: Classification}
%------------------------------------------------------------------------------

\noindent\textbf{Öguz, Socolar, Steinhardt, Torquato (2019).}
This paper provides a comprehensive classification of hyperuniformity and
\textbf{anti-hyperuniformity} in all one-dimensional substitution tilings.
Given a substitution matrix $\mathbf{M}$ with Perron--Frobenius eigenvalue $\lambda_1$
and second-largest eigenvalue $\lambda_2$, the \textbf{eigenvalue formula}
\[
  \alpha = 1 - \frac{2\ln|\lambda_2|}{\ln\lambda_1}
\]
expresses the hyperuniformity exponent entirely in terms of the matrix spectrum.
The exponent $\alpha$ fills $[-1,3]$ densely over the space of all valid substitution
matrices, including anti-hyperuniform cases ($\alpha < 0$) where the structure factor
\emph{diverges} at the origin.

The Pisot--Vijayaraghavan (PV) property --- $|\lambda_i|<1$ for all $i\geq2$ ---
is shown to be necessary and sufficient for the existence of \textbf{Bragg peaks}
(non-zero weight in the diffraction measure at isolated $k$-values).
Non-Pisot substitutions lack Bragg peaks entirely, and their diffraction measure is
purely continuous; $Z(k)$ then decays faster than any power law, placing them
outside the $\alpha$-classification scheme.
Concrete results: Fibonacci $\alpha=3$ (Class~I), Period-Doubling $\alpha=1$ (Class~II).
These two cases are used throughout this project.

\vspace{0.4em}
\noindent\textbf{Bombieri \& Taylor (1986).}
Although no PDF is available, this paper is a key reference: the authors studied the
\emph{cubic} 3-letter substitution $a\to aac$, $b\to ac$, $c\to b$ (equivalently
encoded by a $3\times3$ integer matrix) and showed that it produces a well-defined
diffraction pattern with Bragg peaks.
With $\lambda_1 \approx 2.247$ and $|\lambda_2| \approx 0.802$ the eigenvalue formula
gives $\alpha = 1.545$, the \textbf{first known example filling the gap $1<\alpha<2$}.
This chain was reproduced in the present project with $\bar\Lambda = 0.377$.

%------------------------------------------------------------------------------
\subsection{Quasicrystals and Projection Methods}
%------------------------------------------------------------------------------

\noindent\textbf{Shechtman, Blech, Gratias, Cahn (1984).}
The experimental discovery of an Al-Mn alloy with icosahedral diffraction symmetry ---
forbidden by classical crystallography --- launched the field of quasicrystals.
Bragg peaks appear at irrational wavevector ratios, implying long-range order without
periodicity. This motivates the theoretical framework of \emph{cut-and-project}
(projection) methods in which a higher-dimensional periodic lattice is projected
onto a lower-dimensional physical space along an irrational direction.

\vspace{0.4em}
\noindent\textbf{Öguz, Macé, Torquato (2017).}
This paper analyzes hyperuniformity of projection-method quasicrystals analytically.
For the ideal Fibonacci strip (projection window aligned with the golden ratio),
the structure factor satisfies $Z(k) \sim k^4$, giving $\alpha=3$ (Class~I).
A nonideal strip (irrational window boundary) degrades to $\alpha=1$ (logarithmic variance,
Class~II): the rational approximants to the boundary introduce quasi-periodic modulations
that slow the decay of $S(k)$.

The paper derives an \textbf{analytical upper bound} $\sigma^2(R) \leq 1/4$ for the
Fibonacci chain, consistent with the numerical Fibonacci $\bar\Lambda = 0.201 < 1/4$
found in this project (and with $\bar\Lambda_{\rm Silver} = 1/4$ as a candidate upper-bound
saturation for $n=2$).
The two-point correlation $g_2(r)$ is shown to be piecewise quadratic with inflection
points at the Fibonacci spacings $1$ and $\tau = (1+\sqrt5)/2$.

%------------------------------------------------------------------------------
\subsection{Stealthy Hyperuniform Ground States}
%------------------------------------------------------------------------------

\noindent\textbf{Batten, Stillinger, Torquato (2008).}
This paper introduced the concept of \textbf{collective-coordinate optimization}:
given a target function $\Phi = \sum_{|\mathbf{k}|\leq K} S(\mathbf{k})$, gradient
descent on particle positions drives $S(\mathbf{k}) \to 0$ for $|k| \leq K$,
engineering a \textbf{stealthy} configuration.
The resulting patterns are highly degenerate ground states (many configurations
satisfy the stealth condition) and are generally \emph{disordered}.
Special cases include super-ideal gas variants and equi-luminous structures.
A key insight: stealth can be achieved with particles that appear to have no macroscopic
order, yet exhibit complete isotropy and hidden long-range correlations.
This paper predates the thermodynamic ensemble framework of Torquato et al.\ (2015)
and provides the algorithmic foundation for all stealthy ensemble generation.

\vspace{0.4em}
\noindent\textbf{Torquato, Zhang, Stillinger (2015) --- \emph{Physical Review X}.}
This paper develops the \textbf{exact ensemble theory} for stealthy hyperuniform
disordered ground states.
The control parameter is
\[
  \chi = \frac{M(K)}{d(N-1)},
\]
where $M(K)$ is the number of independently constrained wavevectors within a ball
of radius $K$, $N$ is the particle number, and $d$ the dimension.
For $\chi < 1/2$ the ground-state manifold is \emph{disordered} (positive entropy);
for $\chi > 1/2$ only periodic configurations satisfy the constraints.
All thermodynamic quantities --- isothermal compressibility $\kappa_T$, pair
correlations $g_2(r)$, structure factor $S(k)$ --- are derived analytically in the
ensemble limit.
The paper also derives the relation $K^* = 4\pi\rho\chi$ (for $d=1$) between the
exclusion radius $K^*$ and the control parameter $\chi$.

In $d=1$ the paper's Table~V provides reference values of $\bar\Lambda$ for stealthy
configurations, which are well approximated by $\bar\Lambda \approx 1/(\pi^2\chi)$
--- a formula derived independently in this project and validated against the
collaborator's ensemble data.

%------------------------------------------------------------------------------
\subsection{Perturbed Lattices and Cloaking}
%------------------------------------------------------------------------------

\noindent\textbf{Klatt, Kim, Torquato (2020).}
This paper introduces and analyzes the \textbf{uncorrelated random lattice (URL)}:
each site of the integer lattice $\mathbb{Z}^d$ is displaced independently and
uniformly within a hypercube $[-a/2,a/2]^d$.
In one dimension the \textbf{exact formula}
\[
  \bar\Lambda(a) = \frac{1+a^2}{6}
\]
is derived analytically (for $a \leq 1$) and extended via periodicity to all $a>0$.
Crucially, the URL always has $\alpha=2$ regardless of $a$: the uniform displacement
smears the lattice Bragg peaks into a continuous background, but the origin
$k=0$ retains the same power-law suppression $S(k)\sim k^2$.

The central discovery is \textbf{cloaking}: when the displacement amplitude $a$
equals an integer (1, 2, 3, \ldots), the form factor $\tilde f(k) = \sin(\pi k a)/(\pi k a)$
vanishes at every non-zero reciprocal lattice vector, completely eliminating the
Bragg peaks from the diffraction pattern.
At $a=1$: $\bar\Lambda = 1/3$, which this project identifies as the \emph{exact limit}
of the metallic-mean sequence $\bar\Lambda(\mu_n) \to 1/3$ as $n\to\infty$.
The $\tau$ order metric (based on pair correlations) diverges between cloaking events,
providing a complementary characterization of the structural transition.

%------------------------------------------------------------------------------
\subsection{Diffusion Spreadability}
%------------------------------------------------------------------------------

\noindent\textbf{Torquato (2021) --- \emph{Physical Review E}.}
This paper introduces \textbf{diffusion spreadability} $\mathcal{S}(t)$ as a dynamic
probe of microstructure.
Consider a two-phase medium (phases 1 and 2) with a solute initially confined to
phase 1; $\mathcal{S}(t)$ measures the fraction of solute transferred to phase 2 by
time $t$.
For hyperuniform media with exponent $\alpha$ the long-time behavior is
\[
  E(t) \equiv \mathcal{S}(\infty) - \mathcal{S}(t) \;\sim\; t^{-(d+\alpha)/2},
\]
while stealthy media show \emph{exponential} decay $E(t)\sim e^{-ct}$.
At short times $\mathcal{S}(t) \sim t^{1/2}$ in all cases, providing access to the
\textbf{specific surface} $s$ (surface area per unit volume of the interface).
Applications include NMR/MRI characterization of porous media, where the diffusion
signal encodes structural information across length scales.
The 1D formula $E(t) \sim t^{-(1+\alpha)/2}$ is used in this project to extract $\alpha$
via a log-log regression of the excess spreadability.

\vspace{0.4em}
\noindent\textbf{Wang \& Torquato (2022) --- \emph{Physical Review Applied}.}
This paper applies spreadability to a broad set of 2D and 3D model microstructures:
URL packings, stealthy ground states, Debye random media (nonhyperuniform),
Bravais lattices, and real porous media (Fontainebleau sandstone).
A practical algorithm is presented to extract $\alpha$ from the slope of
$\ln E(t)$ versus $\ln t$ in the asymptotic regime.
For one-dimensional structures the paper gives $\phi_2 = 0.35$ for the Fibonacci
chain (volume fraction of phase~2 in the two-phase embedding), which is used throughout
this project for all 1D spreadability calculations.

\vspace{0.4em}
\noindent\textbf{Hitin-Bialus, Maher, Steinhardt, Torquato (2024).}
This paper achieves high-precision extraction of $\alpha$ for quasiperiodic tilings
via spreadability.
By exploiting the \textbf{self-similarity} of substitution tilings, the small-$k$
Bragg-peak regime (where oscillations dominate) can be truncated, leaving a clean
power-law tail in $E(t)$.
Results: Fibonacci $\alpha = 3.0000 \pm 0.000008$; Period-Doubling
$\alpha = 1.0000 \pm 0.00001$; Penrose tiling (2D) $\alpha = 5.97\pm0.06$
(close to the Class~I value $d+1=3$ for a 2D tiling with the Pisot property).
The paper confirms $\phi_2 = 0.35$ for all 1D chains studied and provides the
most precise numerical validation of the eigenvalue formula~(Öguz et al.\ 2019)
to date.

%------------------------------------------------------------------------------
\subsection{Order Metrics and Comparative Rankings}
%------------------------------------------------------------------------------

\noindent\textbf{Middlemas, Stillinger, Torquato (2019).}
This paper computes $\bar\Lambda$ for all Barlow sphere-packing sequences in $d=3$,
parametrized by fcc/hcp stacking codes.
A linear model $\bar\Lambda(\alpha) = \alpha\bar\Lambda_{\rm fcc} + (1-\alpha)\bar\Lambda_{\rm hcp}$
holds for periodic stacking sequences (where $\alpha$ is the fraction of fcc layers).
The dominant factor in $\bar\Lambda$ is the local cluster geometry: fcc-like
environments ($60^\circ$ tetrahedral) are more hyperuniform than hcp-like
environments ($90^\circ$ tetrahedral).
Theta-series methods are used for all Bragg-peak sums, establishing the
Zachary--Torquato (2009) technique as standard.
The paper demonstrates that $\bar\Lambda$ serves as a sensitive discriminator among
structures that share the same pair-correlation function to short range.

\vspace{0.4em}
\noindent\textbf{Morse, Steinhardt, Torquato (2024).}
This paper studies ordered versus disordered stealthy hyperuniform systems across
spatial dimensions $d = 1$--$4$.
In $d=1$ the densest Fourier-space hard-sphere packing --- which saturates the
stealthy constraint with the fewest degrees of freedom --- corresponds exactly to
a Bravais lattice.
In higher dimensions, disordered stealthy configurations can achieve \emph{lower}
packing density than their crystalline counterparts, counter to the usual expectation.
The \textbf{order metric} $\tau$ (based on pair-distance statistics) is used as a
complement to $\bar\Lambda$: $\tau$ is large for ordered (lattice-like) stealthy
configurations and small for disordered ones.
The results motivate using both $\bar\Lambda$ and $\tau$ to fully characterize
stealthy hyperuniform materials, and demonstrate that stealth alone does not
determine the degree of order.

%------------------------------------------------------------------------------
\subsection{Published Reference Values}
%------------------------------------------------------------------------------

\noindent The table below collects all $\alpha$ and $\bar\Lambda$ values that appear
explicitly in the published literature, restricted to the $d=1$ patterns directly
relevant to this project plus the 2D Penrose tiling.
Values marked \dag{} were confirmed independently by this project.

\vspace{0.4em}
\noindent\begin{tabular}{llccl}
\toprule
\textbf{Pattern} & \textbf{Class} & $\alpha$ \textbf{(published)} & $\bar\Lambda$ \textbf{(published)} & \textbf{Source} \\
\midrule
Integer lattice ($d=1$)
  & I   & $\infty$        & $1/6$ \dag         & Torquato \& Stillinger (2003) \\
Fibonacci chain
  & I   & $3$ (exact) \dag & $0.20110$ \dag    & Zachary \& Torquato (2009); \\
  &     &                 &                    & Öguz et al.\ (2019); Hitin-Bialus et al.\ (2024) \\
Period-Doubling chain
  & II  & $1$ (exact) \dag & diverges          & Öguz et al.\ (2019); Hitin-Bialus et al.\ (2024) \\
Bombieri-Taylor chain
  & I   & $1.545$ \dag    & ---                & Bombieri \& Taylor (1986); Öguz et al.\ (2019) \\
URL ($a > 0$, $d=1$)
  & I   & $2$ \dag        & $(1 + a^2)/6$ \dag & Klatt, Kim, Torquato (2020) \\
URL cloaked ($a = 1$)
  & I   & $2$ \dag        & $1/3$ \dag         & Klatt, Kim, Torquato (2020) \\
Stealthy (d=1, various $\chi$)
  & I   & n/a \dag        & see PRX Table~V    & Torquato, Zhang, Stillinger (2015) \\
Penrose tiling ($d=2$)
  & I   & $5.97 \pm 0.06$ & ---                & Hitin-Bialus et al.\ (2024) \\
\bottomrule
\end{tabular}

\vspace{0.3em}
\noindent\dag{} Independently confirmed by numerical computation in this project.
``---'' indicates the value was not reported in the cited work.
$\bar\Lambda$ values for the Period-Doubling and Bombieri-Taylor chains were
\emph{not} published prior to this project.
The Penrose tiling $\alpha$ is consistent with $d + \alpha_{\rm 1D} = 2 + 3 = 5$ but the
precise interpretation in 2D remains open.

%==============================================================================
\section{Week 1}
%==============================================================================

\subsection{Code Benchmarking and Validation}

The core numerical algorithm for computing $\sigma^2(R)$ was validated against
the 1D Poisson process ($\sigma^2 = 2\rho R$) to within \textbf{1.1\% mean
relative error} across the full $R$-range.
Output: \texttt{fig1\_poisson\_benchmark.png}.

\subsection{Quasicrystal Generation}

Two generation methods were implemented:
\begin{enumerate}[noitemsep]
  \item \textbf{Substitution method} (\texttt{substitution\_tilings.py}): iterative rule
        application up to $N \sim 10^7$ for all five metallic-mean chains.
  \item \textbf{Cut-and-project method} (\texttt{projection\_method.py}): Fibonacci chain
        from $\mathbb{Z}^2$ projection. Agreement with substitution: $<0.1\%$ on $\bar\Lambda$.
\end{enumerate}

\noindent\begin{tabular}{lccc}
\toprule
\textbf{Chain} & \textbf{Metallic mean} $\mu_n$ & $|\det M|$ & $\alpha$ (exact) \\
\midrule
Fibonacci ($n=1$)   & $(1+\sqrt{5})/2 \approx 1.618$ & 1 & 3 \\
Silver    ($n=2$)   & $1+\sqrt{2} \approx 2.414$     & 1 & 3 \\
Bronze    ($n=3$)   & $(3+\sqrt{13})/2 \approx 3.303$& 1 & 3 \\
Copper    ($n=4$)   & $2+\sqrt{5} \approx 4.236$     & 1 & 3 \\
Nickel    ($n=5$)   & $(5+\sqrt{29})/2 \approx 5.193$& 1 & 3 \\
\bottomrule
\end{tabular}

\vspace{0.3em}\noindent All $2\times2$ metallic-mean matrices have $|\det M|=1$,
which forces $|\lambda_2|=1/\lambda_1$ and thus $\alpha=3$ exactly for all of them.

\subsection{Number Variance and $\bar\Lambda$}

$\sigma^2(R)$ was computed for all five metallic-mean chains at $N \sim 10^6$--$10^7$.
All are Class~I (bounded, oscillating variance). The exact formula for the integer
lattice is $\sigma^2(R) = \{R\}(1-2\{R\})$ where $\{R\}$ denotes the fractional part.

\noindent\begin{tabular}{lcc}
\toprule
\textbf{Pattern} & $\bar\Lambda$ & \textbf{Notes} \\
\midrule
Integer Lattice & $1/6 = 0.1667$ & Literature minimum \\
Fibonacci       & $0.2011 \pm 0.001$ & Confirms Zachary \& Torquato (2009) \\
Silver          & $0.2500$ & Novel (later confirmed $= 1/4$ exactly; see §4) \\
Bronze          & $0.2821 \pm 0.0006$ & Novel \\
Copper          & $0.293$ & Novel \\
Nickel          & $0.310$ & Novel \\
\bottomrule
\end{tabular}

\subsection{Spreadability}

Point patterns are embedded into two-phase media ($\phi_2 = 0.35$). The excess
spreadability $E(t) = \mathcal{S}(\infty) - \mathcal{S}(t)$ yields $\alpha$ from
the long-time slope $E(t) \sim t^{-(1+\alpha)/2}$.

\noindent\begin{tabular}{lccc}
\toprule
\textbf{Pattern} & $\alpha$ (measured) & $\alpha$ (expected) & Error \\
\midrule
Fibonacci & 3.049 & 3 & 1.6\% \\
Silver    & 2.992 & 3 & 0.3\% \\
Bronze    & 2.987 & 3 & 0.4\% \\
\bottomrule
\end{tabular}

\vspace{0.3em}\noindent Integer lattice confirmed $\alpha \to \infty$ (exponential decay).
Spreadability was validated as a second, independent measure of $\alpha$.

%==============================================================================
\section{Week 2}
%==============================================================================

\subsection{Uncorrelated Random Lattice and Stable Distributions}

The uncorrelated random lattice (URL) displaces each site by $U(-a,a)$.
The exact formula $\bar\Lambda(a) = (1+a^2)/6$ was derived and validated to $<0.3\%$
for all $a$. URL gives $\alpha=2$ for all $a>0$ (Class~I). At $a=1$ (cloaking limit),
$\bar\Lambda = 1/3$ exactly --- this turns out to be the limit of the metallic-mean sequence.

Symmetric $\alpha$-stable perturbations fill the $1 < \alpha < 2$ space stochastically
(Class~II/III for stability index $s < 1$).

\subsection{Stealthy Hyperuniform Patterns}

Stealthy configurations ($S(k) = 0$ for $|k| < K^*$) were analyzed:
\begin{itemize}[noitemsep]
  \item Collaborator-provided ensemble: $\sim$4300 disordered-phase configurations per
        $\chi$ ($N = 2000$, $\rho = 1$, $\chi \in \{0.1, 0.2, 0.3\}$).
  \item 1D approximation derived: $\bar\Lambda \approx 1/(\pi^2\chi)$, validated below.
\end{itemize}

\noindent\begin{tabular}{cccc}
\toprule
$\chi$ & $\bar\Lambda$ (measured) & $\bar\Lambda = 1/(\pi^2\chi)$ & Excess \\
\midrule
0.10 & 1.021 & 1.013 & $+0.8\%$ \\
0.20 & 0.526 & 0.507 & $+3.7\%$ \\
0.30 & 0.357 & 0.338 & $+5.6\%$ \\
\bottomrule
\end{tabular}

\textbf{Two-phase discovery:} Generating stealthy configurations via collective-coordinate
optimization from a \emph{perturbed lattice} starting point gives $\bar\Lambda \approx 1/6$
--- near the lattice minimum --- while collaborator data (random starting positions) gives
$\bar\Lambda \approx 1/(\pi^2\chi) \gg 1/6$. Both phases satisfy $S(k)=0$ for $|k|<K^*$
but differ drastically in their spatial structure. This shows stealthy hyperuniformity does
not uniquely determine a point pattern.

\textbf{Re-classification:} Stealthy patterns are not power-law suppressed; $\alpha$ is not
meaningful. The correct descriptor is $K^* = 4\pi\rho\chi$.

\subsection{New Substitution Chains}

Three new chains were added to \texttt{substitution\_tilings.py}:

\noindent\begin{tabular}{lllcc}
\toprule
\textbf{Chain} & \textbf{Rules} & \textbf{Matrix} & \textbf{Class} & $\alpha$ \\
\midrule
Period-Doubling & $S \to L$, $L \to LSS$    & $[[0,2],[1,1]]$    & II  & 1.000 \\
0222 Chain      & $S \to LL$, $L \to SSLL$  & $[[0,2],[2,2]]$    & III & 0.639 \\
Bombieri-Taylor & $a \to aac$, $b \to ac$, $c \to b$ & $3\times3$ & I   & 1.545 \\
\bottomrule
\end{tabular}

\textbf{Bombieri-Taylor} fills the $1 < \alpha < 2$ gap:
\[
  M_{BT} = \begin{pmatrix}2&0&1\\1&0&1\\0&1&0\end{pmatrix},\quad
  \lambda_1 = 2.247,\quad |\lambda_2| = 0.802,\quad
  \alpha = 1 - \frac{2\ln(0.802)}{\ln(2.247)} = \mathbf{1.545}, \quad
  \bar\Lambda = 0.377
\]
This is the first known Class~I quasicrystal with $1 < \alpha < 2$.

A new $\alpha=2$ chain was also added (first deterministic substitution tiling at this value,
previously only the stochastic URL occupied $\alpha=2$):
\[
  a \to bc,\; b \to c,\; c \to abc \quad
  \Rightarrow \quad \lambda_1 \approx 2.148,\; |\lambda_{2,3}| \approx 0.682\text{ (complex pair)},\;
  \alpha = 2.000,\; \bar\Lambda = 0.275.
\]

\subsection{Full Catalog Assembly and Junior Paper}

All patterns were compiled into \texttt{catalog.json} via \texttt{research\_catalog.py}.
Infrastructure: \texttt{disordered\_patterns.py}, \texttt{stealthy\_analysis.py},
\texttt{make\_figures.py}.

The Junior Paper (\texttt{jp/jp\_hyperuniformity.tex}) was drafted:
§1 Introduction, §2 Theoretical Background, §3 Methods (5 subsections).
9 pages, biblatex+biber backend. Results sections (§4--§8) exist as a separate draft
(\texttt{jp/jp\_results\_draft.tex}) pending organization.

%==============================================================================
\section{Week 3}
%==============================================================================

Work completed between the second and third meetings.

\subsection{Spreadability Convergence for Bombieri-Taylor}

\textbf{Figure:} \texttt{fig\_bt\_deep.png}

Spreadability converges to $\alpha = 1.545$ for BT, but very slowly:
$\hat\alpha = 2.26$ at $t \in [10^2, 10^5]$, $1.98$ at $[10^4, 10^7]$,
$1.51$ at $[10^8, 10^{11}]$. Two causes:
\begin{enumerate}[noitemsep]
  \item Large $|\lambda_2| = 0.802$ (vs.\ $0.618$ for Fibonacci) means correction
        terms in the spreadability decay slowly.
  \item BT has $24\times$ more Bragg peaks than Fibonacci ($10^5$ vs.\ $4{,}300$),
        causing oscillations that bias slope fits.
\end{enumerate}
$\bar\Lambda$ from $\sigma^2(R)$ avoids these issues (no fitting window needed);
$\bar\Lambda_{\rm BT} = 0.377$.

\subsection{Finite-Size Convergence}

\textbf{Figure:} \texttt{fig\_finite\_size\_alpha.png}

$\bar\Lambda$ from $\sigma^2(R)$ is stable from $N \sim 10^3$ for both BT and Fibonacci.
BT spreadability $\hat\alpha$ converges from ${\sim}2.4$ toward $1.545$ but needs $t > 10^8$.
Fibonacci spreadability works at moderate $t$.

\subsection{$3\times3$ Substitution Matrix Search}

\textbf{Figure:} \texttt{fig\_new\_tilings.png}

Enumerated all 39,304 $3\times3$ non-negative integer matrices with $|\det M| = 1$
and row sums $\leq 4$. Of these, 4,882 satisfy $|\lambda_i| < 1$ for $i \geq 2$.
Among those, 504 have complex conjugate eigenvalue pairs; all 504 give $\alpha = 2$
exactly (18 distinct $\lambda_1$ values, range 1.32--3.85). The remaining 4,378 matrices
(all-real eigenvalues) give $\alpha = 3$. No matrix gives $\alpha \in (2,3)$.

The $\alpha = 2$ result follows from the determinant constraint: for $3\times3$ with
$|\det M| = 1$ and complex pair $\lambda_2 = \bar\lambda_3$,
\[
  \lambda_1|\lambda_2|^2 = 1 \;\Rightarrow\;
  |\lambda_2| = \lambda_1^{-1/2} \;\Rightarrow\; \alpha = 2 \text{ exactly.}
\]

Example: $a \to bc$, $b \to c$, $c \to abc$ ($\lambda_1 \approx 2.148$, $\bar\Lambda = 0.274$).
$\bar\Lambda$ varies widely across the 504 matrices (0.26 to 1.7), even among those
sharing the same $\lambda_1$.

\subsection{Generalized Metrics for Classes II and III}

\textbf{Figure:} \texttt{fig\_generalized\_metrics.png}

$\bar\Lambda$ is not defined for Classes II and III ($\sigma^2(R)$ grows without bound).
Instead:
\begin{align*}
  \Lambda_{II}   &= \lim_{R\to\infty} \frac{\sigma^2(R)}{\ln R} \qquad \text{(Class II)} \\
  \Lambda_{III}  &= \lim_{R\to\infty} \frac{\sigma^2(R)}{R^{1-\alpha}} \qquad \text{(Class III)}
\end{align*}
Extracted via curve fit: $\sigma^2 = \Lambda_{II}\ln R + b$ (Class II) and
$\sigma^2 = \Lambda_{III} R^{1-\alpha} + b$ (Class III, $\alpha$ fixed to theory).
Results: Period-Doubling $\Lambda_{II} = 0.082$;
0222 chain $\Lambda_{III} = 0.143$ (with $\alpha = 0.639$ from eigenvalue formula).

\subsection{The $2 < \alpha < 3$ Gap}

\textbf{Figure:} \texttt{fig\_nonunimodular.png}

Extended the $3\times3$ search to $4\times4$: 24M matrices tested, 1.36M with
$|\lambda_i| < 1$, zero with $\alpha \in (2,3)$. The gap persists because:
\begin{itemize}[noitemsep]
  \item Rank 3 with complex pair: $|\det M| = 1$ forces $\alpha = 2$ (proof above)
  \item Rank $\geq 4$ with complex pair: extra small eigenvalues force $|\lambda_2|$
        larger, giving $\alpha < 2$
  \item All-real eigenvalues: numerical search finds only $\alpha = 3$
\end{itemize}
No substitution matrix with $|\det M| = 1$ at any rank can produce $\alpha \in (2,3)$.

\textbf{Non-unit determinant:} For $2\times2$ matrices with $|\lambda_2| < 1$:
\[
  \alpha = 3 - \frac{2\ln|\det M|}{\ln\lambda_1}
\]
To get $\alpha \in (2,3)$: need $1 < |\det M| < \sqrt{\lambda_1}$. Since det is an
integer, only $|\det M| = 2$ works (requires $\lambda_1 > 4$). Found 28 such matrices,
all with $|\det M| = 2$ and $\alpha \in (2.0, 2.2)$.

Smallest example: $M = \bigl[\begin{smallmatrix}0&1\\2&4\end{smallmatrix}\bigr]$
($\det = -2$), rules $a \to b$, $b \to aabbbb$.
$\lambda_1 = 2+\sqrt{6} \approx 4.449$, $\alpha = 2.071$.
$\sigma^2(R)$ oscillates over $[0.001, 0.87]$ (Fibonacci: $[0.15, 0.25]$).
$\bar\Lambda \approx 0.648 \pm 0.004$ ($N = 17.8$M, period-averaged).
First known 1D Class~I chain with $\alpha \in (2,3)$.

\subsection{Silver $\bar\Lambda = 1/4$ and Metallic-Mean Convergence}

\textbf{Figures:} \texttt{fig\_silver\_lambda\_quarter.png},
\texttt{fig\_metallic\_convergence.png}

\textbf{Silver:} $N = 9.37 \times 10^6$ points, $5 \times 10^4$ windows:
$\bar\Lambda = 0.2502 \pm 0.0006$ ($0.3\sigma$ from $1/4$).
Silver has $\rho = 1/2$ exactly. Fibonacci: $0.201$ (ours) vs.\ $0.20110$
(Zachary \& Torquato 2009, theta-series; both numerical).
Bronze: $0.2822 \pm 0.0006$ (no obvious closed form).

\textbf{Metallic-mean sequence:} Computed $\bar\Lambda(\mu_n)$ for $n = 1$ through $20$:

\noindent\begin{tabular}{cccccccccc}
\toprule
$n$ & 1 & 2 & 3 & 4 & 5 & 6 & 7 & 8 & 10 \\
$\bar\Lambda$ & 0.201 & 0.250 & 0.282 & 0.293 & 0.310 & 0.318 & 0.321 & 0.325 & 0.327 \\
\midrule
$n$ & 15 & 20 & $\infty$ & & & & & & \\
$\bar\Lambda$ & 0.329 & 0.331 & $1/3$ & & & & & & \\
\bottomrule
\end{tabular}

\vspace{0.2em}\noindent Monotone from below. Limit $= 1/3 = \bar\Lambda$ of URL cloaked
lattice ($a = 1$, Klatt et al.\ 2020). As $n \to \infty$, $L/S = \mu_n \to \infty$ but
$\ell_L/\langle\ell\rangle \to 1$: the long tile becomes the mean spacing, short tiles
become rare and tiny, and the chain approaches a lattice with disorder matching URL $a = 1$.

%==============================================================================
\section{Week 4}
%==============================================================================

Tasks from the third meeting (2026-03-27) with Prof.\ Torquato.

\subsection{Why $\alpha_{\max} = 3$ for 1D Substitution Tilings}

The eigenvalue formula $\alpha = 1 - 2\ln|\lambda_2|/\ln\lambda_1$ (Öguz et al.\ 2019)
gives $\alpha = 3$ for all metallic-mean chains. The bound $\alpha \leq 3$ holds
for \emph{all} unimodular ($|\det M| = \pm 1$) substitution matrices, for the following reason.

For a rank-$r$ substitution matrix $M$ with Perron--Frobenius eigenvalue $\lambda_1$
and remaining eigenvalues $\lambda_2, \ldots, \lambda_r$, the determinant constraint gives
\[
  |\det M| = \lambda_1 \prod_{i=2}^{r} |\lambda_i| \geq \lambda_1 |\lambda_2|^{r-1}
  \quad\Rightarrow\quad
  |\lambda_2| \leq \left(\frac{|\det M|}{\lambda_1}\right)^{1/(r-1)}.
\]
Substituting into the eigenvalue formula:
\[
  \alpha = 1 - \frac{2\ln|\lambda_2|}{\ln\lambda_1}
    \leq 1 + \frac{2}{r-1}\left(\frac{\ln\lambda_1 - \ln|\det M|}{\ln\lambda_1}\right)
    = 1 + \frac{2}{r-1}\left(1 - \frac{\ln|\det M|}{\ln\lambda_1}\right).
\]
For $|\det M| = 1$ this simplifies to $\alpha \leq 1 + 2/(r-1)$. The bound is saturated
when all sub-dominant eigenvalues have the same absolute value, i.e.,
$|\lambda_i| = \lambda_1^{-1/(r-1)}$ for $i \geq 2$.

\medskip
\noindent\begin{tabular}{ccl}
\toprule
\textbf{Rank $r$} & $\alpha_{\max}$ & \textbf{Example} \\
\midrule
$2$ & $3$ & all metallic means (Fibonacci, Silver, \ldots) \\
$3$ & $2$ & 504 cubic matrices with complex pair \\
$4$ & $5/3 \approx 1.67$ & (no known examples saturating this) \\
$r \to \infty$ & $\to 1^+$ & approaches Class~II boundary \\
\bottomrule
\end{tabular}

\medskip\noindent
Thus $\alpha = 3$ is the global maximum for unimodular substitution tilings in 1D,
achieved uniquely by rank-2 (two-letter) chains. Higher-rank matrices necessarily
have $\alpha < 3$ because the extra eigenvalues in $|\det M| = 1$ force $|\lambda_2|$
to be larger (closer to 1), reducing $\alpha$. Non-unimodular matrices can exceed $\alpha = 3$
in principle (e.g.\ $|\det M| = 0$ gives $\alpha = \infty$, but this degenerates to a lattice),
though in practice the range $\alpha \in (2,3)$ requires $1 < |\det M| < \sqrt{\lambda_1}$
as shown in §3.5.

\subsection{Physical Interpretation of the $\bar\Lambda$ Ranking}

$\bar\Lambda$ quantifies residual large-scale density fluctuations: larger $\bar\Lambda$
means more disorder. The catalog ranking (§5) is controlled by a few
structural parameters:

\begin{enumerate}[noitemsep]
  \item \textbf{Perfect order $\to$ minimum.} The integer lattice achieves
        $\bar\Lambda = 1/6$, the global minimum for all 1D hyperuniform patterns.
        Any departure from perfect periodicity increases $\bar\Lambda$.

  \item \textbf{Metallic means: tile disparity controls disorder.}
        As $n$ grows, the tile-length ratio $\mu_n = L/S$ increases.
        Short tiles become rarer (frequency $\sim 1/\mu_n \to 0$) and act as
        ``defects'' in an otherwise lattice-like array of long tiles.
        More extreme defects $\to$ larger variance $\to$ larger $\bar\Lambda$.
        The sequence $\bar\Lambda(\mu_1) = 0.201 < \bar\Lambda(\mu_2) = 0.250 < \cdots$
        is monotone increasing toward $1/3$.

  \item \textbf{URL: displacement amplitude $a$ controls disorder directly.}
        Each site is displaced by $U(-a/2, a/2)$, giving
        $\bar\Lambda = (1+a^2)/6$.
        At $a = 0$: lattice ($\bar\Lambda = 1/6$).
        At $a = 1$: cloaked ($\bar\Lambda = 1/3$), matching the metallic-mean limit.
        Intermediate $a$ interpolates smoothly.

  \item \textbf{Stealthy: weaker constraint $\to$ more disorder.}
        The stealth condition $S(k) = 0$ for $|k| < K^*$ constrains fewer
        degrees of freedom when $\chi$ is small.
        $\bar\Lambda \approx 1/(\pi^2\chi)$: smaller $\chi$ allows larger fluctuations.

  \item \textbf{Bombieri-Taylor: slow decay of $S(k)$ means large fluctuations.}
        $\alpha = 1.545$ is low (close to Class~II boundary), so
        $S(k)$ decays slowly near $k = 0$, allowing more long-wavelength
        fluctuations. Combined with $|\lambda_2| = 0.802$ (close to 1),
        giving $\bar\Lambda = 0.377$.
\end{enumerate}

\noindent In summary, $\bar\Lambda$ increases with: (i) larger tile-length disparity
in substitution chains, (ii) larger displacement amplitude in perturbed lattices,
(iii) weaker spectral constraints in stealthy patterns, and (iv) slower power-law
suppression of $S(k)$ near the origin.

\subsection{Uncloaked URL Patterns}

\textbf{Figure:} \texttt{fig\_url\_sweep.png}

The exact formula $\bar\Lambda(a) = (1+a^2)/6$ (Klatt et al.\ 2020) was validated
by simulation at 17 values of $a \in [0.05, 2.0]$ with $N = 50{,}000$ points each.
All simulated values match the formula to within $1\%$.

\medskip
\noindent\begin{tabular}{ccccc}
\toprule
$a$ & $\bar\Lambda$ (exact) & $\bar\Lambda$ (sim) & Error & Note \\
\midrule
0    & 0.1667 & ---    & ---     & lattice \\
0.5  & 0.2083 & 0.2082 & 0.08\% & \\
1.0  & 0.3333 & 0.3330 & 0.10\% & cloaked (Bragg peaks vanish) \\
1.5  & 0.5046 & 0.5047 & 0.02\% & \\
2.0  & 0.6667 & 0.6657 & 0.14\% & second cloaking event \\
\bottomrule
\end{tabular}

\medskip\noindent
Key observations: (i) $\bar\Lambda$ grows monotonically with $a$ (more displacement
$=$ more disorder). (ii) At each integer $a = 1, 2, \ldots$, cloaking occurs: the
Bragg peaks of the underlying lattice vanish because the displacement form factor
$\tilde f(k) = \sin(\pi k a)/(\pi k a)$ has zeros at all non-zero reciprocal lattice
vectors. (iii) $\bar\Lambda(a=1) = 1/3$ matches the metallic-mean limit exactly.

\subsection{$g_2$-Invariant Patterns}

Torquato \& Stillinger (2003, §V.B) define a \textbf{$g_2$-invariant process}
as one where the pair correlation function $g_2(r)$ remains invariant as
density increases, up to a terminal density $\rho_c$ at which the system
becomes hyperuniform ($S(0) = 0$).

Three 1D examples were investigated:

\begin{enumerate}[noitemsep]
  \item \textbf{Lattice-gas:} Tessellate $[0, L)$ into $N$ intervals of length $D = 1$;
        place one point uniformly at random in each interval.
        $g_2(r) = r/D$ for $r \leq D$, $g_2 = 1$ for $r > D$.
        This is equivalent to URL with $a = 1$ (just a phase offset:
        $j + U(0,1)$ vs.\ $j + U(-1/2, 1/2)$).
        \textbf{$\bar\Lambda = 0.333 = 1/3$}, confirming the URL result.

  \item \textbf{Step-function $g_2$:} $g_2(r) = \Theta(r - D)$ at terminal density
        $\phi_c = 1/2^d$. For $d = 1$: $\phi_c = 1/2$, close-packed hard rods $=$ lattice.
        \textbf{$\bar\Lambda = 1/6$}.

  \item \textbf{Gaussian shuffled lattice:} Each lattice site displaced by
        $\mathcal{N}(0, \sigma^2)$. Hyperuniform for any $\sigma > 0$.
        Measured: $\bar\Lambda = 0.200$ ($\sigma = 0.1$),
        $0.250$ ($\sigma = 0.2$), $0.342$ ($\sigma = 0.3$).
        At $\sigma = 0.1$, $\bar\Lambda$ matches Fibonacci;
        at $\sigma = 0.2$, it matches Silver $= 1/4$.
\end{enumerate}

\noindent The $g_2$-invariant processes at terminal density are all hyperuniform
and fall naturally into the existing ranking: they are particular cases of
perturbed lattices, with $\bar\Lambda$ determined by the displacement distribution.
The lattice-gas is a special case of URL. The step+delta $g_2$ examples
(Eq.~121--125 in the 2003 paper) are 3D constructions with no direct 1D analog
beyond the cases already computed.

\subsection{Metallic Chain Tile Visualizations}

\textbf{Figure:} \texttt{fig\_metallic\_tiles\_extended.png}

Figure shows 25 tiles of each metallic-mean chain for $n = 1, 2, 3, 4, 5, 6, 8, 10$.
As $n$ increases, short tiles (S) become rarer: the S fraction decreases from $36\%$
at $n = 1$ (Fibonacci) to $8\%$ at $n = 10$.
For large $n$, the chain approaches a near-lattice of long tiles with rare,
widely-spaced short-tile defects.

\subsection{Metallic-Mean Convergence to $n = 100$}

Extended $\bar\Lambda(\mu_n)$ from $n = 20$ to $n = 100$ (TARGET\_N $= 300$k).

\medskip
\noindent\begin{tabular}{cccccccccccc}
\toprule
$n$ & 1 & 2 & 3 & 5 & 10 & 20 & 30 & 50 & 70 & 100 & $\infty$ \\
$\bar\Lambda$ & 0.199 & 0.250 & 0.283 & 0.310 & 0.327 & 0.332 & 0.332 & 0.334 & 0.330 & 0.335 & $1/3$ \\
\bottomrule
\end{tabular}

\medskip\noindent
All values are consistent with $1/3$ within statistical error at large $n$.
Some nominal overshoots at $n \geq 50$ are within the $\sim\!0.005$ error bars
(reduced accuracy at large $n$ because chains grow very slowly:
$\lambda_1 \approx n + 1/n$, so even at $n = 100$ only two substitution iterations
are needed to reach $N = 300$k, limiting the quasi-periodic structure).
The overall picture is unchanged: \textbf{$\bar\Lambda(\mu_n) \to 1/3$
monotonically from below}, confirming the URL cloaked-lattice limit.

\subsection{All Non-Unimodular Chains}

\textbf{Script:} \texttt{compute\_all\_nonunimodular.py}

Enumerated all $2\times2$ non-negative integer matrices with $|\det M| = 2$,
$|\lambda_2| < 1$, $\alpha \in (2,3)$, and row sums $\leq 10$.
Found \textbf{138 matrices} (more than the original 28 found with row sums $\leq 6$)
across 12 distinct $\alpha$ values.

\medskip
\noindent\begin{tabular}{cccc}
\toprule
$\alpha$ & Count & $\bar\Lambda$ range & Smallest $\bar\Lambda$ \\
\midrule
2.071 & 12 & $[0.32, 2.01]$ & $M = \bigl[\begin{smallmatrix}0&2\\1&9\end{smallmatrix}\bigr]$ \\
2.087 & 10 & $[0.34, 1.57]$ & \\
2.175 & 18 & $[0.33, 3.36]$ & \\
2.199 & 14 & $[0.34, 3.57]$ & \\
2.248 & 10 & $[0.33, 2.82]$ & \\
2.271 & 16 & $[0.34, 3.76]$ & \\
2.301 & 18 & $[0.33, 4.05]$ & \\
2.343 & 12 & $[0.33, 4.79]$ & \\
2.362 & 12 & $[0.34, 4.80]$ & \\
2.376 &  8 & $[0.33, 4.97]$ & \\
2.393 &  4 & $[0.34, 5.08]$ & \\
\bottomrule
\end{tabular}

\medskip\noindent
Key observations: (i) All $\alpha$ values lie in $[2.07, 2.39]$, consistent with
the formula $\alpha = 3 - 2\ln 2/\ln\lambda_1$ and the constraint $\lambda_1 > 4$.
(ii) $\bar\Lambda$ varies enormously \emph{within} each $\alpha$ group (by factors of 5--15),
showing that $\alpha$ alone does not determine $\bar\Lambda$.
(iii) The smallest $\bar\Lambda \approx 0.32$--$0.34$ across all groups, close to $1/3$
(the URL cloaked value).
(iv) These are all Class~I ($|\lambda_2| < 1$, bounded $\sigma^2$).

\subsection{Yuan-Torquato Improved Spreadability}

\textbf{Script:} \texttt{yuan\_improved\_spreadability.py}

Implemented the improved fitting procedure from Yuan \& Torquato (2026, arXiv:2602.17873).
Three fitting function types were tested on Fibonacci ($\alpha = 3$):

\medskip
\noindent\begin{tabular}{llccc}
\toprule
\textbf{Chain} & \textbf{Method} & $\hat\alpha$ & $|\hat\alpha - \alpha|$ & Error \\
\midrule
\multirow{3}{*}{Fibonacci ($\alpha = 3$)}
  & Simple log-slope      & 2.19  & 0.81 & 27\% \\
  & Type~I, order 1       & 3.04  & 0.04 & 1.3\% \\
  & Type~III, order 1     & 2.95  & 0.05 & 1.5\% \\
\midrule
\multirow{3}{*}{BT ($\alpha = 1.545$)}
  & Simple log-slope      & 2.27  & 0.72 & 47\% \\
  & Type~I, order 6       & 1.57  & 0.03 & \textbf{1.7\%} \\
  & Type~III, order 6     & 2.16  & 0.61 & 40\% \\
\bottomrule
\end{tabular}

\medskip\noindent
For Fibonacci, Type~I at low order ($n = 1$) gives the best result ($1.3\%$ error);
higher orders diverge because quasicrystal $s^{ex}(t)$ oscillates (dense Bragg peaks).
For BT, Type~I at high order ($n = 5$--$6$) dramatically reduces the error from $47\%$
to $\sim\!2\%$ --- the \textbf{first time BT's $\alpha$ has been measured accurately from
spreadability without extreme time ranges} ($t > 10^8$ was previously needed).
Type~III fails for BT because the spectral density is not analytic at the origin.

\textbf{Conclusion:} The Yuan-Torquato Type~I fitting is a significant improvement
for chains with large $|\lambda_2|$ (like BT), where correction terms dominate at
moderate~$t$. It does not replace variance-based $\bar\Lambda$ extraction but provides
a viable alternative for $\alpha$ measurement from spreadability data.

%==============================================================================
\section{Complete Pattern Catalog}
%==============================================================================

Patterns sorted by $\bar\Lambda$ within Class~I.

\vspace{0.4em}
\noindent\begin{tabular}{llcc}
\toprule
\textbf{Pattern} & \textbf{Class} & $\boldsymbol\alpha$ & $\boldsymbol{\bar\Lambda}$ \\
\midrule
\rowcolor{resultbg}
Integer Lattice             & I & $\infty$ & $1/6 = 0.167$ \\
URL ($a = 0.1$)             & I & 2.0 & 0.168 \\
URL ($a = 0.3$)             & I & 2.0 & 0.181 \\
Gaussian $\sigma=0.1$       & I & 2.0 & 0.187 \\
Fibonacci                   & I & 3.0 & 0.201 \\
URL ($a = 0.5$)             & I & 2.0 & 0.208 \\
\rowcolor{resultbg}
Silver                      & I & 3.0 & $\mathbf{1/4} = 0.250$ \\
\rowcolor{newbg}
Cubic $\alpha$=2 (new)      & I & 2.0 & 0.275 \\
Bronze                      & I & 3.0 & 0.282 \\
Copper                      & I & 3.0 & 0.293 \\
Nickel                      & I & 3.0 & 0.310 \\
URL ($a = 1.0$)             & I & 2.0 & 0.333 \\
Stealthy ($\chi = 0.3$)     & I & ---$^\dagger$ & 0.358 \\
\rowcolor{resultbg}
Bombieri-Taylor             & I & 1.545 & 0.377 \\
Stable $s=1.7$              & I & 1.7  & 0.242 \\
Stable $s=1.5$              & I & 1.5  & 0.303 \\
Stable $s=1.3$              & I & 1.3  & 0.417 \\
Stealthy ($\chi = 0.2$)     & I & ---$^\dagger$ & 0.526 \\
Stealthy ($\chi = 0.1$)     & I & ---$^\dagger$ & 1.021 \\
\midrule
Period-Doubling             & II  & 1.0  & div.\ ($\Lambda_{II} = 0.080$) \\
\midrule
0222 Chain                  & III & 0.639& div.\ ($\Lambda_{III} = 0.209$) \\
Cauchy $\gamma=0.1$         & II  & 0.95 & div. \\
Stable $s=0.7$              & III & 0.7  & div. \\
Stable $s=0.5$              & III & 0.5  & div. \\
Stable $s=0.3$              & III & 0.3  & div. \\
\bottomrule
\end{tabular}

\vspace{0.4em}
\noindent$^\dagger$ Stealthy: $S(k) = 0$ by construction; $\alpha$ not applicable.
Green rows: literature-confirmed benchmarks. Yellow: novel additions.

%==============================================================================
\section{Figures and Scripts}
%==============================================================================

\subsection*{Figures}

\noindent\begin{tabular}{ll}
\toprule
\textbf{Filename} & \textbf{Content} \\
\midrule
\texttt{fig1\_poisson\_benchmark.png}   & Poisson variance validation \\
\texttt{fig2\_bounded\_variance.png}    & $\sigma^2(R)$ for Fibonacci/Silver/Bronze \\
\texttt{fig3\_hyperuniformity\_test.png}& Normalized variance across classes \\
\texttt{fig9\_stealthy\_sk.png}         & $S(k)$ for stealthy at $\chi=0.1,0.2,0.3$ \\
\texttt{stealthy\_Sk\_overlay.png}      & $S(k)$ overlay of all stealthy configs \\
\texttt{bombieri\_taylor\_analysis.png} & BT variance, $S(k)$, spreadability \\
\texttt{fig12\_ranking.png}             & Full $\bar\Lambda$ ranking, Class~I \\
\texttt{fig\_variance\_II\_III.png}     & $\sigma^2(R)$ for Period-Doubling and 0222 \\
\midrule
\texttt{fig\_bt\_deep.png}              & BT oscillations, Bragg peak density, $E(t)$ \\
\texttt{fig\_finite\_size\_alpha.png}   & $\bar\Lambda(N)$, $\Lambda_{II}(N)$, $\hat\alpha(N)$ (portrait) \\
\texttt{fig\_finite\_size\_alpha\_landscape.png} & Same, landscape $1\times4$ for slides \\
\texttt{fig\_new\_tilings.png}          & Cubic $\alpha$=2 chain variance and $S(k)$ \\
\texttt{fig\_generalized\_metrics.png}  & Normalized variance plateaus (Class II/III) \\
\texttt{fig\_finite\_size\_lambda.png}  & Lattice $\bar\Lambda$ analytic vs.\ simulated \\
\texttt{fig\_catalog\_updated.png}      & Updated $\bar\Lambda$ ranking with stealthy fixes \\
\texttt{fig\_metallic\_convergence.png} & $\bar\Lambda(\mu_n)$ for $n=1\ldots20$ converging to $1/3$ \\
\texttt{fig\_4x4\_alpha\_dist.png}      & $4\times4$ matrix search: $\alpha$ distribution \\
\texttt{fig\_nonunimodular.png}         & $\sigma^2(R)$ and running $\bar\Lambda$ for $\alpha=2.071$ chain \\
\bottomrule
\end{tabular}

\subsection*{Scripts}

\noindent\begin{tabular}{ll}
\toprule
\textbf{Script} & \textbf{Purpose} \\
\midrule
\texttt{substitution\_tilings.py}     & All chains (metallic means, BT, cubic $\alpha$=2, non-unimodular) \\
\texttt{quasicrystal\_variance.py}    & $\sigma^2(R)$, $\bar\Lambda$, generalized metrics \\
\texttt{projection\_method.py}        & Cut-and-project Fibonacci chain \\
\texttt{two\_phase\_media.py}         & Spectral density, spreadability \\
\texttt{disordered\_patterns.py}      & URL, stealthy generation, exact $\bar\Lambda$ for URL \\
\texttt{stealthy\_analysis.py}        & Collaborator stealthy ensemble analysis \\
\texttt{research\_catalog.py}         & Full pipeline $\to$ \texttt{catalog.json} \\
\texttt{make\_figures.py}             & Figures from catalog + stealthy results \\
\midrule
\texttt{bt\_deep\_dive.py}            & (Mtg 2--3) BT oscillations, $S(k)$ density, $E(t)$ \\
\texttt{finite\_size\_scaling.py}     & (Mtg 2--3) $\bar\Lambda(N)$, $\Lambda_{II}(N)$, $\hat\alpha(N)$ \\
\texttt{search\_cubic\_tilings.py}    & (Mtg 2--3) Exhaustive $3\times3$ matrix search \\
\texttt{generalized\_metrics.py}      & (Mtg 2--3) $\Lambda_{II}$, $\Lambda_{III}$; patches variance module \\
\texttt{finite\_size\_lambda\_bar.py} & (Mtg 2--3) Lattice formula; stealthy two phases \\
\texttt{metallic\_mean\_convergence.py} & (Mtg 2--3) $\bar\Lambda(\mu_n)$ for $n=1\ldots20$ \\
\texttt{search\_4x4\_tilings.py}      & (Mtg 2--3) $4\times4$ matrix search (24M matrices) \\
\texttt{silver\_precision\_analysis.py} & (Mtg 2--3) High-precision Silver/Bronze $\bar\Lambda$ \\
\texttt{fig\_nonunimodular\_chain.py} & (Mtg 2--3) Figure for $\alpha=2.071$ non-unimodular chain \\
\bottomrule
\end{tabular}

\subsection*{Presentations and Written Work}

\noindent\begin{tabular}{lll}
\toprule
\textbf{Document} & \textbf{Location} & \textbf{Notes} \\
\midrule
Week 1 presentation & \texttt{results/week1/week1\_presentation.tex} & Beamer, core metallic-mean results \\
Week 2 presentation & \texttt{results/week2/week2\_presentation.tex} & Beamer, 12 slides \\
Week 3 presentation & \texttt{results/week3/week3\_presentation.tex} & Beamer, 16 slides (rewritten 2026-03-24) \\
Junior Paper        & \texttt{jp/jp\_hyperuniformity.tex}           & 9 pages, §1--§3 complete \\
JP results draft    & \texttt{jp/jp\_results\_draft.tex}            & §4--§8 + appendix, held for later \\
\bottomrule
\end{tabular}

\vspace{0.3em}\noindent All documents compile via MiKTeX
(\texttt{C:/Users/minec/AppData/Local/Programs/MiKTeX/miktex/bin/x64/}).

%==============================================================================
\section{Key Theoretical Results}
%==============================================================================

\begin{enumerate}
  \item \textbf{$\alpha = 3$ universality of metallic means.} For any $2\times2$
        metallic-mean matrix with $|\det M| = 1$: $|\lambda_2| = 1/\lambda_1 \Rightarrow
        \alpha = 3$ universally.

  \item \textbf{Bombieri-Taylor fills $1 < \alpha < 2$.} The 3-letter cubic substitution
        $a \to aac, b \to ac, c \to b$ gives $\alpha = 1.545$, the first known 1D
        quasicrystal with $\alpha \in (1,2)$.

  \item \textbf{First deterministic $\alpha = 2$ substitution tiling.} The chain
        $c \to abc, a \to bc, b \to c$ gives $\alpha = 2.000$ exactly,
        $\bar\Lambda = 0.275$.

  \item \textbf{$\bar\Lambda$ density invariance.} $\bar\Lambda$ is invariant under
        uniform density rescaling in $d=1$; direct comparison across patterns with
        different densities is valid.

  \item \textbf{Integer lattice has no finite-$N$ correction.} $\bar\Lambda = 1/6$
        exactly for all integer $R_\text{max}$ (piecewise cubic integral identity).

  \item \textbf{Stealthy two-phase structure.} Both ordered and disordered stealthy
        phases satisfy $S(k) = 0$ for $|k| < K^*$, but differ drastically in $\bar\Lambda$
        ($\approx 1/6$ vs.\ $1/(\pi^2\chi)$). Stealthy $\alpha$ is not a power-law exponent.

  \item \textbf{$2 < \alpha < 3$ gap is inaccessible to unimodular matrices.}
        Exhaustive search over $3\times3$ (39,304) and $4\times4$ (24M) matrices
        confirms zero unimodular cases with $\alpha \in (2,3)$.
        Rank-3 proof: $|\det M|=1 \Rightarrow \alpha=2$ exactly for complex pairs.
        Rank-$\geq4$ proof: additional Pisot factor forces $\alpha < 2$.

  \item \textbf{General 2$\times$2 formula.} For any $2\times2$ Pisot matrix:
        $\alpha = 3 - 2\ln|\det M|/\ln\lambda_1$. Gap $(2,3)$ is accessible when
        $|\det M| \in (1, \sqrt{\lambda_1})$; all 28 known examples have $|\det M|=2$.

  \item \textbf{First 1D quasicrystal with $\alpha \in (2,3)$.}
        $M = \bigl[\begin{smallmatrix}0&1\\2&4\end{smallmatrix}\bigr]$ ($\det=-2$)
        gives $\alpha = 2.071$, $\bar\Lambda \approx 0.650$ (rough estimate).
        Large log-periodic oscillations in $\sigma^2(R)$, period $\ln\lambda_1 \approx 1.493$.

  \item \textbf{Silver $\bar\Lambda = 1/4$ exact (numerical).}
        $N = 3.88\times10^6$, $3\times10^4$ windows: $\bar\Lambda = 0.2501 \pm 0.0004$
        ($0.3\sigma$ from $1/4$). Analytical proof unknown.

  \item \textbf{$\bar\Lambda(\mu_n) \to 1/3$ confirmed.} Metallic-mean sequence
        $n=1\ldots20$ converges to $1/3$ at rate $0.22/n^{1.43}$. URL cloaked
        lattice ($a=1$, $\bar\Lambda=1/3$) is the exact limit.
\end{enumerate}

%==============================================================================
\section{Open Questions}
%==============================================================================

\begin{enumerate}
  \item \textbf{$2 < \alpha < 3$ gap: \emph{RESOLVED} (mechanism and first example).}
        Non-unimodular $M=\bigl[\begin{smallmatrix}0&1\\2&4\end{smallmatrix}\bigr]$
        gives $\alpha=2.071$.
        Sub-question: what is the minimal-complexity chain in $(2,3)$?
        (Row sum $\geq 6$ is currently the minimum requirement.)

  \item \textbf{Silver exact value: \emph{RESOLVED numerically}.}
        $\bar\Lambda = 1/4$ at $0.3\sigma$. Analytical proof not yet known.

  \item \textbf{$\bar\Lambda$ metallic-mean limit: \emph{RESOLVED}.}
        $\bar\Lambda(\mu_n) \to 1/3$; rate $0.22/n^{1.43}$; URL cloaked lattice
        is the exact limit.

  \item \textbf{Non-unimodular $\bar\Lambda$:} Is $\bar\Lambda = 0.650$ exactly $13/20$?
        Analytical proof unknown; large oscillation amplitude makes high-precision
        measurement difficult.

  \item \textbf{Bronze closed form:} Is $\bar\Lambda(\text{Bronze}) = 0.2822$ a clean
        expression? Closest simple fraction: $11/39 = 0.28205$ ($0.18\sigma$, non-obvious).

  \item \textbf{Stealthy two-phase literature:} Is the ordered/disordered bifurcation
        discussed in existing stealthy ground-state literature?

  \item \textbf{2D generalization:} Does the eigenvalue formula
        $\alpha = 1 - 2\ln|\lambda_2|/\ln\lambda_1$ extend to 2D substitution
        tilings (Penrose, Ammann-Beenker)?
\end{enumerate}

\end{document}
